% FDUletter_example.tex - an example latex file to illustrate FDUletter.cls
%
% Template by Brian Wood (brian.wood@oregonstate.edu).  Please feel free to send suggestions for changes; this template/cls is not exactly elegantly done!
% Modified by Huang Weiran (huangweiran1998@outlook.com) to fit the need of FDU students.
% Modified by Phosphorus15 (phosphophate@gmail.com) into one-page for CSU students.


\documentclass[12pt]{LGUletter}
\usepackage{fontspec} 
\usepackage{tikz} 
\usepackage{xcolor}
	\definecolor{lgupurple}{rgb}{0.5, 0.1, 0.5}
\usepackage{lipsum}
\usepackage{fancyhdr}
\usepackage{lastpage}
\usepackage{eso-pic}
%
% This section is just a bunch of busywork so that the second and following pages read ``Page X of Y''
\pagestyle{fancy}
\fancyhf{}
\renewcommand{\headrulewidth}{0pt}
\renewcommand{\footrulewidth}{0pt}
\rhead{Page \thepage \hspace{1pt} of \pageref{LastPage}}
%
%
% Set custom font here. Comment this line out if you do not have a Cambria font (originally included with this template) installed; computer modern (or whatever your current default font is) will be substituted.
%
%\setmainfont{[Cambria.ttf]}[BoldFont  = [CambriaBold.ttf], ItalicFont  = [CambriaItalic.ttf], BoldItalicFont = [CambriaBoldItalic.ttf] ]

\newcommand{\watermark}[3]{\AddToShipoutPictureBG{
\parbox[b][\paperheight]{\paperwidth}{
\vfill%
\centering%
\begin{tikzpicture}
    \path (0,0) -- (\paperwidth,\paperheight);
    \node[opacity=.08] at (current page.center)
    {\includegraphics[width=0.55\textwidth]{background.pdf}};
    \end{tikzpicture}
\vfill}}}

% The material below is a whole big dang thing whose purpose is just to set up a fixed coordinate system for \tikz so that you can put the Department or School address in the upper right-hand side without it moving all around every time you change something in the page.  I think it works.
% Defining a new coordinate system for the page:
%
% --------------------------
% |(-1,1)    (0,1)    (1,1)|
% |                        |
% |(-1,0)    (0,0)    (1,0)|
% |                        |
% |(-1,-1)   (0,-1)  (1,-1)|
% --------------------------
\makeatletter
\def\parsecomma#1,#2\endparsecomma{\def\page@x{#1}\def\page@y{#2}}
\tikzdeclarecoordinatesystem{page}{
    \parsecomma#1\endparsecomma
    \pgfpointanchor{current page}{north east}
    % Save the upper right corner
    \pgf@xc=\pgf@x%
    \pgf@yc=\pgf@y%
    % save the lower left corner
    \pgfpointanchor{current page}{south west}
    \pgf@xb=\pgf@x%
    \pgf@yb=\pgf@y%
    % Transform to the correct placement
    \pgfmathparse{(\pgf@xc-\pgf@xb)/2.*\page@x+(\pgf@xc+\pgf@xb)/2.}
    \expandafter\pgf@x\expandafter=\pgfmathresult pt
    \pgfmathparse{(\pgf@yc-\pgf@yb)/2.*\page@y+(\pgf@yc+\pgf@yb)/2.}
    \expandafter\pgf@y\expandafter=\pgfmathresult pt
}
\makeatother
%
%
%%%%%%%%%%% Put Personal Information Here %%%%%%%%%%%
%
\def\name{Jiafeng Ding,\\
School of Electronic Information, Central South University
}
%
% List your degree(s), licences, etc. here if you like.
%\def\What{, Your degrees, etc.} 
%
% Set the name of your Department or School here
% I honestly don't know why the negative spacing is necessary to get the alignment of the first line correct.  This must be a ``\tikz'' thing.
%%%%%%%%%%%%%%%%%%  School or Department %%%%%%%%%%%%%%%
\def\Where{\hspace{-18.2mm}\textbf{\color{lgupurple}%{lgupurple}
School of Science and Engineering
}} 
%%%%%%%%%%%%  Additional Contact Information %%%%%%%%%%%
%
% Set your preferred primary contact address here.
\def\Address{Yaohui Han,} 
%
\def\CityZip{Research Assistant} 
%
% Set your FDU e-mail here
\def\Email{\textbf{\color{lgupurple}E-mail}: han0129@csu.edu.cn}
%
%
% Set your department or personal website here
\def\URL{\textbf{\color{lgupurple}URL}: yaohui.netlify.app}
%
%%%%%%%%%%%%%%%%  Footer information  %%%%%%%%%%%%%%%%%%
%
%  The next line is for your college, used as a footer.  If you prefer not to have this, just comment out these lines in favor of the line labeled "[[Alternate]]" below
\def\school{\small{} } 
% \def\school{~}  % [[Alternate]]
%
%%%%%%%%%%%%%%%%%%%%%  Signature line  %%%%%%%%%%%%%%%%%%%%%
%
% Set your signature line here.
% One can add a signature image in a PDF file using the following code; this requires a file called "signature_block.pdf" to be installed in the same folder as the .tex file.  The vertical spacing (\vspace) and the scaling will have to be adjusted to get things to look correct for your particular signature image. Alternatively, comment out the following line in favor of the one labeled "[[Alternate]]" if you want to sign a paper copy of the letter.
%
\signature{ 
\vspace{-12mm}\includegraphics[scale=0.08]{signature_block.pdf}\\
\vspace{2mm}
\name}
%\signature{\name}  % [[Alternate]]

% This block sets up the address on the right-hand side of the header. 
%
% The following lines just compile the information you set up into the LaTex letter variable "address" for later use.
%
%The following command "clears out" the default address so that it can be better set using \tikz
\address{}

\def\newaddress{
\Where\\ 
\Address\\ 
\CityZip\\ 
\Email\\ 
2001 Longxiang Boulevard,\\ Shenzhen, Guangdong
}
%
%%%%%%%%%%%  DATE  %%%%%%%%%%%%%%%%%%%%%%%%%
% If you want a date different from the current date, comment out the next line in favor of the line labeled "[[Alternate]]".  The ``\vspace{10mm}'' just moves the date down a tiny bit so it doesn't interfere with the header.  This can be adjusted to your preference.
%
% \date{\today}
%\date{\vspace{10mm} 20 September 2020}  %[[Alternate]]
%
%%%%%%%%%%% Set the subject here if there is one  %%%%
%\subject{Stuff} % optional subject line

\begin{document}
%
%
%%%%%%%%  The "To" address goes here.
%
\begin{letter}{
               Some University\\ 
               Some Addresss\\ 
               SomeTown, SomeState 					       				  		 ~~SomeZip
               }
% This line sets up the return address to the right-side of the OSU logo.  The location is set with absolute node addresses using ``\tikz''.  It can still be a bit fussy, and you may need to alter this a little to get things to look right.  The bit that changes the position are the numbers in parentheses ``at (14.2,2.7)''
%
\begin{tikzpicture}[remember picture,overlay,,every node/.style={anchor=center}]
\node[text width=7cm] at (page cs:0.5,0.73){\small \newaddress};
\end{tikzpicture} 

%%%%%%  The ``opening'' is just the method of address you would like to use at the start of the letter.
%
\opening{Dear Admissions Committee,}

%%%%%%%%%% Body of letter   %%%%%%%%%%%%%%
% Remove it if you do not want watermark
\watermark{}{}{}

% The ``\lipsum[1-8]" command just fills the letter with 8 paragraphs of Latin for the purposes of filler.  Unless you really want to send filler Latin to someone, you will replace this command with actual text.  Do that here:
%
It is my pleasure to recommend Jane Doe for admission to [Name of Program] at [University]. I am an associate professor of [Subject] at [University]. I came to know Jane while teaching her in [Course]. The course comprised [short description of course]. Jane distinguished herself by submitting an exceptionally well researched and interesting project on [project]. She is one of the very best students that I have taught in the past five years in respect of her writing ability and research skills. 

Overall, Jane is highly intelligent and has good problem-solving skills. Her [project] demonstrated her ability to come a detailed understanding of [details about coursework] the lifestyles of pre-historic hominids, and to analyze the consequences of those practices on evolution. She gave a particularly interesting discussion of the difficult practice of tool-shaping and a profound engagement with the use and the implications of these tools. Her overall intelligence is also reflected in her grades for the course, which were among the highest in the class. 

Jane has excellent communication skills. Her written work is both clear and concise, as well as interesting to read. She demonstrated her oral articulateness in the discussion sections that were an integral part of the course. Each discussion section focused on a particular aspect of Olduvian Gorge hominids. Students were required analyze unsolved questions and to develop and defend their own theories. Jane was highly proficient in applying the course material in analyzing these “puzzles.” She always explained her views very concisely and gave supporting arguments that were both clear and persuasive. Jane also demonstrated good teamwork skills in group assignments. 

At a personal level, Jane is a well-disciplined, industrious student with a pleasant personality. Throughout the course, Jane demonstrated great perseverance and initiative. Not only was she interested in and motivated to learn the material, but she also put great work into assimilating it to her own experience and developing her own ideas about each topic that we discussed. 

If I can be of any further assistance, or provide you with any further information, please do not hesitate to contact me. 

%%%%%%% ``closing'' sets the sign-off line.
\closing{Sincerely,}

% Comment out/in the lines below as necessary
%\encl{If an enclosure is provided, let them know what it is.}

% \ps{A postscript if that is a thing you do.}

% \cc{Someone Who Cares (and is copied).}

\end{letter}

\end{document}